% Part: counterfactuals
% Chapter: introduction
% Section: material-conditional

\documentclass[../../../include/open-logic-section]{subfiles}

\begin{document}

\olfileid{cnt}{int}{mat}

\olsection{مادی شرطیہ}

انگریزی وچ اپنی سبھ توں سادی شکل وچ شرطیہ ایہو جیہا جملہ اے:
’جے \dots تاں \dots‘، جتھے \dots{} آپ جملے ہوندے نیں۔ مثلاً
’جے ایہ کم گھر دے وڈے نوکر نے کیتا اے، تاں مالی بے قصور اے۔‘ منطق
دے ابتدائی کورساں وچ اسیں $ \lif$ رابطہ ورت کے شرطیے علامتاں وچ لکھنے
سکھدے آں۔ \dots نال دسّے حصیاں نوں، مثلاً !!{formula}s $!A$ تے~$!B$ نال
لکھو، تے پورے شرطیے نوں $!A \lif !B$ نال لکھدے آں۔

رابطہ $\lif$ \emph{صداقتی-فنکشنی} اے، یعنی سچائی دی قدر، $\True$ یا $\False$،
$!A \lif !B$ لئی $!A$ تے~$!B$ دیاں سچائی قدراں نال طے ہوندی اے۔
$!A \lif !B$ اُدوں تے صرف اُدوں سچ اے جدوں $!A$ جھوٹ یا $!B$ سچ اے،
تے ہور صورت وچ جھوٹ اے۔ سچائی دی قدر بندی~$\pAssign v$ دے لحاظ نال
اسیں ایویں تعریف کردے آں کہ $\pSat{v}{!A \lif !B}$ اُدوں تے صرف اُدوں اے جدوں
$\pSat/{v}{!A}$ یا $\pSat{v}{!B}$۔ ایس معنیات والے رابطے $\lif$ نوں
\emph{مادی شرطیہ} آکھدے آں۔

ایس تعریف توں منطق دیاں کئی بنیادی گلاں نکلدیاں نیں۔ سبھ توں پہلاں، استنباط دا
قضیہ مادی شرطیے لئی قائم اے:
\begin{align}
  \text{ جے } \Gamma, !A \Entails !B \text{ تاں } \Gamma \Entails !A \lif !B
\end{align}
ایہ صداقتی-فنکشنی اے۔ $!A \lif !B$ تے $\lnot !A \lor !B$ ہم ارز نیں:
\begin{align}
  !A \lif !B & \Entails \lnot !A \lor !B\\
  \lnot !A \lor !B & \Entails !A \lif !B
  \intertext{مادی شرطیہ اپنی تالی دا وی منطقی نتیجہ اے تے اپنے مقدم دی نفی دا وی:}
  !B & \Entails !A \lif !B\\
  \lnot !A & \Entails !A \lif !B
  \intertext{جھوٹا مادی شرطیہ اپنے مقدم تے اپنی تالی دی نفی دے عطف دے ہم ارز اے۔ جے $!A \lif !B$ جھوٹ اے، تاں $!A \land \lnot !B$ سچ اے، تے ایہدا اُلٹ وی:}
  \lnot(!A \lif !B) & \Entails !A \land \lnot !B\\
  !A \land \lnot !B & \Entails \lnot(!A \lif !B)
  \intertext{مادی شرطیے لئی اثباتِ مقدم دا قاعدہ چلدا اے:}
  !A, !A \lif !B & \Entails !B
  \intertext{مادی شرطیاں دیاں تالیاں اکٹھیاں ہو سکدیاں نیں:}
  !A \lif !B,  !A \lif !C & \Entails !A \lif (!B \land !C)
  \intertext{اسیں مقدم نوں ہمیشہ ہور مضبوط کر سکدے آں، یعنی شرطیہ \emph{یک رُخا} اے:}
  !A \lif !B & \Entails (!A \land !C) \lif !B
  \intertext{مادی شرطیہ متعدی اے، یعنی لڑی والا قاعدہ معتبر اے:}
  !A \lif !B, !B \lif !C & \Entails !A \lif !C
  \intertext{مادی شرطیہ اپنے عکسِ نقیض دے ہم ارز اے:}
  !A \lif !B & \Entails \lnot !B \lif \lnot !A\\
  \lnot !B \lif \lnot !A & \Entails !A \lif !B
\end{align}

ایہ سارے استنباط ریاضیاتی دلیل وچ کم آؤندے نیں تے کوئی اوکڑ نہیں پاندے۔
پر فلسفے تے لسانیات دیاں لکھتاں وچ غیر ریاضیاتی موقعیاں اُتے اوہناں دے برابر
استنباطاں دیاں کئی مبینہ ردّی مثالاں دتیاں گئیاں نیں۔ اوہ ایہ اشارہ دیندیاں نیں
کہ انگریزی دے ’جے \dots تاں \dots‘ والے بیاناں نوں علامتاں وچ لکھن لئی
مادی شرطیہ~$\lif$ مناسب رابطہ نہیں، یا گھٹ توں گھٹ ہمیشہ مناسب نہیں۔

\end{document}
